\section{Mejoras Arquitectónicas e Implementación del Modelo Híbrido}
\label{sec:implementation}

El diseño de un Small Language Model basado en una arquitectura híbrida requiere consideraciones cuidadosas sobre cómo integrar efectivamente dos paradigmas computacionales distintos: la atención global de los Transformers y la memoria secuencial eficiente de las unidades recurrentes. Este capítulo documenta el proceso iterativo de desarrollo: qué limitaciones presentaba la propuesta inicial y qué decisiones de ingeniería se tomaron para superarlas. La arquitectura resultante de este proceso se describe formalmente en el Capítulo~\ref{sec:arquitectura}.

\subsection{Arquitectura Original y sus Limitaciones}

La propuesta inicial de una arquitectura híbrida combinaba Transformers y GRU en cada bloque, pero presentaba una falla fundamental: los estados ocultos del GRU se descartaban al final de cada capa, eliminando la capacidad de memoria recurrente real que justificaba su inclusión. En el flujo original:

\begin{enumerate}
    \item El GRU procesaba la entrada y generaba un estado oculto
    \item Este estado oculto se descartaba inmediatamente
    \item En el siguiente bloque, el GRU volvía a inicializar desde cero
    \item Como resultado, no había información compartida verticalmente (entre bloques) a través del componente recurrente
\end{enumerate}

Además, el GRU procesaba \textit{después} de la atención, causando redundancia: la atención ya había capturado dependencias globales, y el GRU intentaba refinar patrones secuenciales sobre una representación ya modificada. Finalmente, sin control explícito de gradientes, el modelo sufría de explosiones típicas de arquitecturas recurrentes durante el entrenamiento temprano.

\subsection{Mejora 1: Persistent GRU Hidden States (CRÍTICA)}

La primera y más importante mejora es hacer que los estados ocultos del GRU sean \textbf{persistentes} a través de todos los 12 bloques del modelo. En lugar de descartar el hidden state al final de cada bloque, ahora se mantiene un tensor de hidden states (uno por bloque) que fluye a través de la red durante todo el forward pass.

\subsubsection{Implementación Técnica}

En la clase \texttt{HybridBlock}, el método \texttt{forward} fue reescrito para retornar tanto la salida como el nuevo hidden state. Cada bloque mantiene conexiones residuales que aseguran que tanto la información de atención como la del GRU se preservan y combinan:

\begin{itemize}
    \item \textbf{Pre-LayerNorm}: Normalización aplicada antes de cada sub-componente para estabilidad numérica
    \item \textbf{GRU Processing}: Captura patrones secuenciales locales, mantiene estructura sintáctica
    \item \textbf{Attention Refining}: Refina con contexto global, integrando información de todo el contexto
    \item \textbf{MLP Refinement}: Procesamiento de características no-lineal final con factor de expansión 4
\end{itemize}

En la clase \texttt{HybridModel}, el forward pass ahora mantiene un registro de estos estados a través de todos los bloques. La inicialización de estos estados es crítica: cuando se comienza el entrenamiento, todos los hidden states se inicializan como None, permitiendo que el GRU comience sin sesgo previo. Sin embargo, después del primer forward pass, estos estados contienen información aprendida que se propaga al siguiente paso.

\subsubsection{Impacto en la Memoria Recurrente}

Con esta mejora, cada bloque GRU ahora \textbf{recuerda} información del bloque anterior. Si el Bloque 1 identifica un patrón secuencial importante (ej: estructura sintáctica de una frase), ese patrón se representa en su hidden state y se propaga al Bloque 2. El Bloque 2 puede entonces usar esa información para refinar su comprensión local, creando una verdadera arquitectura recurrente vertical.

Esto es conceptualmente diferente de una RNN pura, donde la información fluye solo horizontalmente (a través de la secuencia). Aquí, la información fluye tanto horizontalmente (dentro de cada bloque) como verticalmente (entre bloques a través de GRU), proporcionando una forma de ``corto circuito recurrente'' que complementa la atención. Este flujo de información vertical es lo que permite que una arquitectura con solo 12 capas (versus 30 en SmolLM2) pueda mantener poder expresivo comparable.

\subsection{Mejora 2: Reorden de Componentes (GRU Pre-Attention)}

La segunda mejora arquitectónica cambia el orden de procesamiento dentro de cada bloque. El diseño inicial situaba la MHA en primer lugar, lo que respondía a una intuición razonable: aprovechar el paralelismo de la atención para capturar contexto global y dejar al GRU el refinamiento secuencial posterior, combinando eficiencia paralela con capacidad recurrente.

\subsubsection*{Por qué la intuición inicial era incorrecta}

La MHA no es un paso de extracción pasiva: es una transformación profunda que redistribuye la representación de cada token en función de todo el contexto, disolviendo la estructura local del input. El GRU, situado después, recibe una señal en la que las dependencias locales —orden de palabras, concordancia morfológica, estructura sintáctica inmediata— ya han sido absorbidas y mezcladas globalmente. Sin información local no procesada, el GRU no puede aportar nada cualitativamente diferente a lo que la MHA ya ha capturado. Ambos componentes compiten por representar lo mismo en lugar de especializarse, generando \textbf{redundancia} y desaprovechando la capacidad del GRU.

Existe además un segundo problema: los estados ocultos del GRU en el diseño original se descartaban al final de cada bloque (descrito en la Mejora 1), lo que eliminaba cualquier posibilidad de memoria vertical recurrente independientemente del orden de los componentes.

\subsubsection*{La solución: flujo local $\to$ global}

Invertir el orden resuelve ambos problemas estructuralmente. El GRU, al operar sobre la representación sin transformar, se especializa en exactamente aquello para lo que está diseñado: capturar dependencias secuenciales locales con coste $O(L)$ y propagar su estado oculto al bloque siguiente. La MHA recibe entonces una representación ya estructurada localmente y puede dedicar toda su capacidad a resolver dependencias globales a larga distancia. El resultado es una división de responsabilidades explícita:

\begin{enumerate}
    \item \textbf{GRU} (primero): captura patrones secuenciales locales —orden de palabras, concordancia morfológica, estructura sintáctica—, operando sobre el estado oculto persistente del bloque anterior.
    \item \textbf{MHA} (segundo): integra dependencias globales a larga distancia sobre una representación ya parcialmente refinada localmente, sin necesidad de modelar también la sintaxis inmediata.
    \item \textbf{MLP} (tercero): proyección no-lineal final con factor de expansión 4 ($768 \rightarrow 3072 \rightarrow 768$).
\end{enumerate}

Este orden refleja mejor cómo se estructura el procesamiento lingüístico: primero relaciones locales (concordancia, orden), luego semántica global (co-referencia anafórica, relaciones de largo alcance). La información fluye de lo local a lo global dentro de cada bloque, y verticalmente entre bloques a través del estado oculto del GRU.

\subsection{Mejora 3: Inicialización Ortogonal para Componentes RNN}

Las unidades recurrentes son notoriamente sensibles a la inicialización de pesos. Una inicialización inadecuada puede llevar a vanishing gradients o exploding gradients durante el backpropagation a través del tiempo (BPTT).

La mejora implementa inicialización ortogonal para los pesos del GRU, mientras mantiene inicialización normal estándar para componentes no-recurrentes. Específicamente:

\begin{itemize}
    \item Pesos \texttt{weight\_ih} (input-to-hidden): inicializados como matriz ortogonal
    \item Pesos \texttt{weight\_hh} (hidden-to-hidden): inicializados como matriz ortogonal
    \item Bias: inicializados a cero
    \item Todos los demás módulos (Linear, Embedding, LayerNorm): inicialización estándar
\end{itemize}

La inicialización ortogonal preserva la magnitud de los vectores a través de las transformaciones lineales del GRU, lo que reduce el riesgo de vanishing/exploding gradients. Las matrices ortogonales tienen la propiedad de que $Q^T Q = I$, lo que significa que la transformación preserva normas y ángulos. Para redes recurrentes, esto es crítico porque permite que gradientes fluyan sin atenuación o amplificación artificial a través de múltiples pasos de tiempo.

\subsection{Mejora 4: Gradient Clipping (CRÍTICA)}

Aunque la inicialización ortogonal ayuda, las redes recurrentes siguen siendo propensas a explosiones de gradiente. Por ello, se implementó un mecanismo de clipping de gradientes que limita la norma de los gradientes durante el entrenamiento.

El clipping se aplica después de cada backward pass y antes del optimizer step, asegurando que ningún gradiente explosion pueda causar actualizaciones de parámetros descontroladas. El parámetro \texttt{max\_norm} (típicamente 1.0) define el umbral máximo. Si la norma de gradientes excede este valor, todos los gradientes se escalan proporcionalmente hacia abajo.

Matemáticamente, si $\|\nabla L\| > \text{max\_norm}$, entonces $\nabla L \leftarrow \nabla L \cdot \frac{\text{max\_norm}}{\|\nabla L\|}$.

Este método es especialmente importante en arquitecturas híbridas donde el GRU puede amplificar gradientes a través de sus múltiples iteraciones internas. Sin clipping, un paso de tiempo con gradientes grandes puede causar que todos los 12 bloques GRU amplífiquen ese gradiente, llevando a explosión. Con clipping, esta amplificación es limitada.

\subsection{Mejora 5: Logging de Contribuciones de Componentes}

Para entender cómo cada componente (GRU, Attention, MLP) contribuye al entrenamiento, se implementó un sistema de logging que rastrea la norma de gradientes para cada uno en tiempo real.

Durante cada paso de entrenamiento, después del backward pass y clipping, se calcula:

\begin{itemize}
    \item \textbf{GRU norm}: $\sqrt{\sum_{\text{params in GRU}} (\nabla param)^2}$
    \item \textbf{Attention norm}: $\sqrt{\sum_{\text{params in Attention}} (\nabla param)^2}$
    \item \textbf{MLP norm}: $\sqrt{\sum_{\text{params in MLP}} (\nabla param)^2}$
    \item \textbf{GRU/Attention ratio}: $\frac{\text{GRU norm}}{\text{Attention norm} + \epsilon}$
\end{itemize}

El ratio GRU/Attention es la métrica clave. Permite validar que la arquitectura es verdaderamente híbrida:

\begin{itemize}
    \item \textbf{Ratio < 0.3}: GRU demasiado débil, indica que la atención domina completamente, el GRU no está aprendiendo
    \item \textbf{Ratio 0.5-1.5}: \textbf{BALANCEADO}, ambos componentes contribuyen significativamente al aprendizaje
    \item \textbf{Ratio > 1.5}: GRU dominando, indica desbalance donde GRU absorbe la mayoría del aprendizaje
\end{itemize}

El logging se produce cada N pasos configurables (típicamente 100), permitiendo monitoreo en tiempo real. Un log típico se ve como:

\texttt{[Step 100] Loss: 4.2341 | GRU: 0.4521 | Attn: 0.3841 | MLP: 0.2154 | Ratio: 1.176 | GradNorm: 2.341->1.000}

\subsection{Mejora 6: Validation Loop}

Para detectar overfitting temprano y validar que el modelo está aprendiendo efectivamente, se implementó una evaluación periódica en un conjunto de validación durante el entrenamiento.

La validation loop se ejecuta cada N pasos (típicamente 500), evaluando el modelo en un mini-batch de validación. Esto permite:

\begin{itemize}
    \item Detectar divergencia entre training loss y validation loss (indicador de overfitting)
    \item Validar que el entrenamiento es estable
    \item Tomar decisiones sobre cuándo detener el entrenamiento
\end{itemize}

Se espera que durante el entrenamiento correcto, tanto training loss como validation loss desciendan. Si validation loss comienza a aumentar mientras training loss sigue bajando, indica overfitting en los datos de entrenamiento.

\subsection{Mejora 7: Herramienta de Análisis Arquitectónico}

Se creó la clase \texttt{HybridArchitectureAnalyzer} que proporciona capacidades de diagnóstico detalladas más allá del logging básico. Esta herramienta permite:

\begin{itemize}
    \item Análisis completo de contribuciones de componentes
    \item Exportación de métricas a JSON para análisis posterior con herramientas externas
    \item Generación de reportes legibles por humanos
    \item Análisis del flujo de hidden states a través de capas
    \item Comparación de arquitecturas alternativas
\end{itemize}

El analyzer puede ser invocado durante o después del entrenamiento, permitiendo tanto monitoreo en vivo como análisis post-hoc de los datos de entrenamiento.

\subsection{Mejora 8: Suite de Tests Automatizados}

Para garantizar la corrección de la implementación, se creó una suite de 7 tests automatizados que validan cada aspecto de la arquitectura:

\begin{enumerate}
    \item \textbf{Forward Pass}: Verifica que el output tiene las dimensiones correctas
    \item \textbf{Loss Computation}: Valida que la pérdida se calcula correctamente
    \item \textbf{Gradient Flow}: Asegura que gradientes fluyen correctamente a través de todos los componentes
    \item \textbf{Hidden State Persistence}: Verifica que los hidden states se mantienen a través de bloques
    \item \textbf{Weight Initialization}: Confirma que los pesos se inicializan correctamente
    \item \textbf{Component Analysis}: Valida que el analyzer funciona correctamente
    \item \textbf{Model Size}: Confirma que el conteo de parámetros es el esperado
\end{enumerate}


\subsection{Impacto Cuantificable en la Convergencia}

Las mejoras arquitectónicas tienen impacto medible y significativo en la convergencia del modelo durante el entrenamiento desde cero.

\subsubsection{Velocidad de Convergencia}

\begin{itemize}
    \item \textbf{Arquitectura Original}: Loss 5.4 → 5.1 → 4.8 (primeros 300 steps)
    \item \textbf{Arquitectura Mejorada}: Loss 5.4 → 4.9 → 4.3 (primeros 300 steps)
    \item \textbf{Mejora}: ~15\% más rápido en convergencia inicial
\end{itemize}

Esta mejora en velocidad es significativa porque acelera todo el ciclo de entrenamiento. Con 5 épocas de entrenamiento, una mejora del 15\% en convergencia temprana se traduce en horas ahorradas.

\subsubsection{Estabilidad de Gradientes}

\begin{itemize}
    \item \textbf{Arquitectura Original}: Gradient spikes erráticos, máximo 45+, completamente sin control
    \item \textbf{Arquitectura Mejorada}: Controlados mediante clipping en rango [0.5, 1.0]
    \item \textbf{Mejora}: 45× más estable, entrenamiento predecible y reproducible
\end{itemize}

La estabilidad es crítica para evitar colapso del entrenamiento. Sin ella, el modelo puede alcanzar estados no-recuperables donde los parámetros divergen.

\subsubsection{Visibilidad y Debuggabilidad}

\begin{itemize}
    \item \textbf{Arquitectura Original}: Solo visible loss total, imposible saber si componentes funcionan correctamente
    \item \textbf{Arquitectura Mejorada}: GRU norm, Attention norm, MLP norm, ratio cada 100 steps, completamente transparente
    \item \textbf{Mejora}: Arquitectura totalmente transparente y debuggeable
\end{itemize}

\subsection{Configuración Recomendada para Entrenamiento}

Basado en las mejoras implementadas, la configuración recomendada para entrenar el modelo es:

\begin{verbatim}
from app.training.trainer import TrainingService

service = TrainingService(dataset_path=".datasets/mixed_dataset")

service.train(
    epochs=5,
    batch_size=64,
    learning_rate=1e-4,
    grad_clip_norm=1.0,          # Control de gradientes
    log_metrics_every=100,       # Log de métricas
    validate_every=500,          # Validación periódica
)
\end{verbatim}

Con esta configuración:

\begin{itemize}
    \item Los gradientes se mantienen controlados mediante clipping
    \item Las métricas de cada componente se registran regularmente para monitoreo
    \item La validación detecta overfitting temprano
    \item Los hidden states del GRU persisten y contribuyen significativamente al aprendizaje
\end{itemize}

Las mejoras descritas en este capítulo consolidan el proceso de desarrollo de la arquitectura. El diseño final resultante, con sus especificaciones técnicas completas y la estrategia de entrenamiento, se detalla en el capítulo siguiente.